% 現行研究資料：旧中心商特異性論文VI 第21--22節からの回収
% 分類結果の固有の証明と外部入力（とりわけ有限中心・Azumaya・正則性）を区別する。
\section{三つの特異点の射影特異圏}
\label{sec:quotient-singcat}

本節では$A_\alpha$を中心三次商の一混合項族
\[
 A_\alpha=\frac{\CC[u,v]\langle x,y,z\rangle}
 {([x,y]-z^2,\,[y,z]-(x^2+\alpha uv),\,[z,x]-y^2,\,
   \Omega_1+\alpha uvx+u^3+v^3)}
\]
とする。ただし
\[
 \Omega_1=\frac{x^3+y^3+z^3}{3}
       -\frac{xyz+xzy+yxz+yzx+zxy+zyx}{9}
\]
であり、$u,v$は一次中心元である。
相対PBW性と周囲のAS正則性は
系\ref{cor:quotient-ambient-pbw}で証明した。
中心有限性・Azumaya locusと二重被覆の比較は
命題\ref{prop:quotient-azumaya}、\ref{prop:chart-descent}で証明した。
以下で用いる境界の正則性とIwanaga--Gorenstein性も、
次の補題でこれらの入力から導く。

\begin{lemma}[境界の正則性とchartのGorenstein性]
\label{lem:quotient-boundary-gorenstein}
$\alpha\ne0$とする（比較構成は命題\ref{prop:chart-descent}の末尾の一般化を用いる）。
$A_\alpha/(u)$のqgrは大域次元2を持ち、
$u$-chartの$t=v/u=0$での商$S_1/(\Omega_1+1)$は
大域次元高々2を持つ。
また$Q_\alpha=A_\alpha[u^{-1},v^{-1}]_0$は
左右noetherianで中心上有限なIwanaga--Gorenstein代数であり、
左右自己入射次元は高々3である。
\end{lemma}
\begin{proof}
$A_\alpha/(u)=S_1[v]/(\Omega_1+v^3)$であり、
$v$-chartは$S_1/(\Omega_1+1)$、境界は$S_1/(\Omega_1)$である。
Euler導分$\delta(\Omega_1)=3\Omega_1$と
補題\ref{lem:gluing-moving-central-value}でchartの大域次元は高々2。
境界は滑らかな楕円曲線のcoherent圏なので、
命題\ref{prop:gluing-one-chart}と
補題\ref{lem:gluing-serre-dimension}からqgrの大域次元は2となる。

命題\ref{prop:chart-descent}の階数2の中心\'etale被覆で
\[
 Q'_\alpha\simeq
 R_1[s,s^{-1}]/\left(g+\tfrac32(s^{-3}+s^3/\alpha^3)\right).
\]
$R_1[s,s^{-1}]$は左右noetherianで大域次元高々4。
括弧内は非零中心正則元なので、中心正則元に対するReesの補題は
任意の有限生成左右$Q'_\alpha$加群$M$について
\[
 \operatorname{Ext}^i_{Q'_\alpha}(M,Q'_\alpha)
 \simeq
 \operatorname{Ext}^{i+1}_{R_1[s,s^{-1}]}
                  (M,R_1[s,s^{-1}])
\]
を与える。従って左右自己入射次元は高々3。
中心平坦基底変更で有限生成加群のExtは可換であり、
忠実平坦性がその消滅を検出するので、この上界は$Q_\alpha$へ降下する。
中心上有限性・noetherian性は命題\ref{prop:chart-descent}による。
\end{proof}

\newcommand{\SingCat}{\operatorname{D}_{\mathrm{sg}}}
\newcommand{\PerfCat}{\operatorname{Perf}}
\newcommand{\thickcat}{\operatorname{thick}}
\newcommand{\Kar}{\mathrm{Kar}}
\newcommand{\Db}{\mathrm D^b}
$\mathcal P_\alpha=\thickcat\{\pi A_\alpha(n):n\in\mathbb Z\}
\subset\Db(\operatorname{qgr}A_\alpha)$と置き、
\[
 \mathcal S_\alpha=
 \Db(\operatorname{qgr}A_\alpha)/\mathcal P_\alpha
\]
を射影的特異圏とする。上付き$\Kar$はKaroubi完備化を表す。

\begin{lemma}[中心元に沿う特異圏の局所化]
\label{lem:quotient-central-localization}
$H$は左右noetherianな次数付き代数、
$z\in Z(H)_1$は正則元とする。
$\Db(\operatorname{qgr}(H/(z)))$が構造対象の全twistで
厚く生成されるなら、中心局所化は
\[
 \left(\Db(\operatorname{qgr}H)/
       \thickcat\{\pi H(n)\}\right)^\Kar
   \simeq \SingCat(H[z^{-1}]_0)^\Kar
\]
を与える。次数なしの左右noetherian代数$\Lambda$と
中心正則元$z$に対し、$\Lambda/(z)$の大域次元が有限なら
同じ結論が成立する。
\end{lemma}
\begin{proof}
中心局所化で零となる有限生成加群は
$z$のある冪で殺される。$z$の冪による有限filtrationの
各商は$H/(z)$上のcoherent加群なので、仮定と
$0\to H(-1)\xrightarrow{z}H\to H/(z)\to0$
によって局所化の導来核は$\thickcat\{\pi H(n)\}$に入る。
局所化後の有限生成加群には有限個の生成元の
$H$-spanというlatticeが存在する。
有界複体の微分と射・homotopyの有限個の分母も
$z$の冪で払えるので、有界導来圏の局所化は
$z$-torsion cohomologyを零にするVerdier局所化となる。
twistsは局所化後に自由加群となり、逆に有限自由複体は
twistsの有限和とconeから得られる。
局所化後のperfect対象はこれらの直和因子なので、
直和因子の持上げを仮定せず両辺をKaroubi完備化すればよい。
非次数付きの場合は有限次元torsionによる商を省く。
\end{proof}

\begin{corollary}[二つの中心chartの除去]
\label{cor:quotient-two-chart}
補題\ref{lem:quotient-boundary-gorenstein}の下で
\[
 \mathcal S_\alpha^\Kar\simeq\SingCat(Q_\alpha)^\Kar .
\]
\end{corollary}
\begin{proof}
$A_\alpha/(u)\simeq
 S_1[v]/(\Omega_1+v^3)$は
$v$の非零スカラー変換で既知の正則曲面のモデルに同型である。
同補題の大域次元二とtwistsの厚い生成性により、
まず補題\ref{lem:quotient-central-localization}を$z=u$に適用する。
$u$-chartにおける$t=0$の商は
$S_1/(\Omega_1+1)$であり、
Euler導分$\delta(\Omega_1)=3\Omega_1$を用いる
補題\ref{lem:gluing-moving-central-value}により有限大域次元を持つ。
同補題のaffine版を$t$に適用して$Q_\alpha$に至る。
共通境界$u=v=0$も最初のfiltrationに含まれる。
\end{proof}

\begin{lemma}[孤立行列nodeの完備化と直和分解]
\label{lem:quotient-completion}
$C$をaffine可換noetherian環、
$\Lambda$を$C$上有限な左右noetherian Iwanaga--Gorenstein代数とする。
有限個の極大イデアル$\mathfrak m_i$を除いて
$\Lambda_{\mathfrak p}$が有限大域次元であり、
各完備化が$\widehat\Lambda_{\mathfrak m_i}\simeq
\operatorname{Mat}_{d_i}(N_i)$で、
$N_i$は孤立通常二重点の完備局所環とする。
さらに各局所nodeの単純剰余加群を与える
$\Lambda/\mathfrak m_i\Lambda$の有限加群が存在するとする。
このとき
\[
 \SingCat(\Lambda)^\Kar\simeq
 \bigoplus_i\SingCat(N_i)^\Kar .
\]
\end{lemma}
\begin{proof}
Iwanaga--Gorenstein性に対するBuchweitz--Happel--Chenの同値
\cite{ChenRelativeSingcat}により、
左辺は有限生成Gorenstein-projective加群のstable圏で計算できる。
$M,N$がそのような加群ならstable Homは
\[
 \underline{\operatorname{Hom}}_\Lambda(M,N)=
 \operatorname{coker}\!\left(
 N\otimes_\Lambda\operatorname{Hom}_\Lambda(M,\Lambda)
 \longrightarrow\operatorname{Hom}_\Lambda(M,N)\right)
\]
で、有限$C$加群である。
例外点以外では$M$はGorenstein-projectiveかつ
有限射影次元なので射影的である。
従ってstable Homは有限個の閉点だけに支持されて有限長。
有限表示性、中心の平坦完備化、有限長加群の局所分解により、
各Hom群は各完備化のstable Homの直和に同型となる。
同じ議論はshiftにも適用でき、完備化は完全忠実である。

本質的全射性をKaroubi完備化後に証明する。
通常二重点$N_i=k[[z_1,\ldots,z_e]]/(q)$の
maximal Cohen--Macaulay加群を行列因子化で表すと、
$d^2=q$の微分から$\partial_{z_j}q$は特異圏で零作用となる。
非退化二次式ならこれは$z_j$を零化するので、
$z_j$のKoszul複体は元の対象を直和因子として含み、
そのcohomologyは剰余体の有限直和である。
従ってnodeの特異圏は剰余体で厚く生成される。
仮定した有限$\Lambda$加群はその生成元を完備化の像として与え、
直和因子まで閉じれば所望の圏全体になる。
\end{proof}

\begin{theorem}[三つの通常二重点に集中する射影特異圏]
\label{thm:quotient-three-singcat}
$28\alpha^3+243=0$とし、
上記のPBW・既知Sklyanin入力を用いる。
中心局所化の前提とIwanaga--Gorenstein性は
補題\ref{lem:quotient-boundary-gorenstein}、
三つの局所模型とそれ以外の点での有限大域次元は
定理\ref{thm:three-homological}から従う。
このとき
\[
 \mathcal S_\alpha^\Kar\simeq
 \bigoplus_{\tau^3=1}
  \SingCat\!\left(
   \frac{\CC[[a,b,c,d]]}{a^2+b^2+c^2+d^2}
  \right)^\Kar
 \simeq \PerfCat_{\mathbb Z/2}(\CC)^{\oplus3}.
\]
従って$K_0(\mathcal S_\alpha^\Kar)\simeq\mathbb Z^3$で、
非零直既約対象は同型を除いて六つである。
$[1]$は各二つの直既約対象を交換し、$[2]\simeq\mathrm{id}$。
\end{theorem}
\begin{proof}
系\ref{cor:quotient-two-chart}で$Q_\alpha$に局所化する。
局所模型\ref{thm:three-homological}と
補題\ref{lem:quotient-completion}から最初の圏同値を得る。
ここでは中心上有限な代数の剰余環が$\Mat_6(\CC)$なので、
そのcolumn単純加群を$Q_\alpha$加群と見れば、
補題\ref{lem:quotient-completion}の生成元の存在条件も満たされる。

二次元node $N_0=k[[p,q]]/(pq)$のMCM加群は
正規化$k[[p]]\oplus k[[q]]$上の飽和を取り、
conductorと剰余体の線形部分空間を標準化することにより、
直和因子として$N_0$, $k[[p]]$, $k[[q]]$からなる。
stable圏では$N_0$は零、後の二つは相互直交で
stable Endが$\CC$、shiftは二つを交換する。
よって$\SingCat(N_0)^\Kar\simeq\PerfCat_{\mathbb Z/2}(\CC)$。
まずaffineな$k[p,q]/(pq)$と$k[p,q,r,s]/(pq+rs)$に
OrlovのKn\"orrer周期性\cite[Theorem~2.1]{OrlovSingcat}を適用する。
両方の特異点は原点だけなので、補題\ref{lem:quotient-completion}を
可換環に適用して完備化へ移せる。従って
$N_0$と$k[[p,q,r,s]]/(pq+rs)$の
Karoubi完備な特異圏は同値である。
$\CC$上で後者の二次式は四変数の平方和と線形同型なので結論。
二周期的な一因子の$K_0$は$\mathbb Z$で、
三因子の直和が残りの主張を与える。
\end{proof}
\begin{remark}
全ての同値はKaroubi完備化\emph{後}の主張である。
$D^b(\operatorname{qgr}A_\alpha)$そのものや
非可換cubicのKuznetsov成分と同値であるとは主張しない。
中心の21個の特異点と、代数のホモロジー的特異点三個を区別する。
\end{remark}

\section{正則三次商の球面ツイストと構造対象の標識}
\label{sec:quotient-spherical-twist}
本節は$\alpha=1$の正則例を対象とする。
$A=A_1$、$\mathcal O=\pi A$とし、
正則性$\operatorname{gldim}\operatorname{qgr}A=3$は
定理\ref{thm:quotient-exact-tails}から従う。
命題\ref{prop:quotient-ambient-nakayama}と
補題\ref{lem:gluing-serre-dimension}の双対性から
$\mathbb S_{\Db(\operatorname{qgr}A)}\simeq(-2)[3]$も従う。
さらにfiltered三変数代数
\[
 R_1=
 \frac{\CC\langle X,Y,Z\rangle}
 {([X,Y]-Z^2,\ [Y,Z]-(X^2+1),\ [Z,X]-Y^2)},
 \quad G_1=\Omega_1(X,Y,Z)+X
\]
に、命題\ref{prop:local-two-axial}から得られる三次元単純左加群の線形双対である\emph{右}加群$V$を取る。この加群は
\[
 \dim_\CC V=3,\qquad VG_1=0,\qquad
 \operatorname{End}_{R_1}(V)=\CC,\quad
 \operatorname{Ext}^1_{R_1}(V,V)=0
\tag{QS}
\]
を満たす。単純性と$VG_1=0$は命題\ref{prop:local-two-axial}、無限小剛性は命題\ref{prop:quotient-axial-rigidity}で証明される。後者の法$29$の厳密小行列式は現在の検算コードから再現できる。左加群から右加群への双対は有限次元自己拡大と単純性を保つ。

$t=v/u$とし、
$Q^\circ=A[u^{-1},v^{-1},(v^3-u^3)^{-1}]_0$と置く。
$t=-1$、$(\bar x,\bar y,\bar z)=(iX,iY,iZ)$
（$i^2=-1$）により$V$を$Q^\circ$-加群とする。
次数付き右$A$-加群$N_V$を
\[
 (N_V)_n=
 \begin{cases}V,&n\ge0,\\0,&n<0,\end{cases}
 \qquad u=\mathrm{id},\ v=-\mathrm{id},\
 x=iX,\ y=iY,\ z=iZ
\]
という次数一作用で定め、$\mathcal T=\pi N_V$と書く。
$\mathcal T$は非零coherent対象で、
$u$がtailで可逆なので$\mathcal T(1)\simeq\mathcal T$。

\begin{lemma}[局所自己拡大の消滅]
\label{lem:quotient-rigid-chart}
(QS)を用いると
\[
 \operatorname{End}_{Q^\circ}(V)=\CC,\qquad
 \operatorname{Ext}^1_{Q^\circ}(V,V)=0 .
\]
さらに命題\ref{prop:quotient-azumaya}と
命題\ref{prop:chart-descent}の$\alpha=1$での比較から
$\operatorname{gldim}Q^\circ\le3$。
\end{lemma}
\begin{proof}
自己拡大$0\to V\to E\to V\to0$に対し、
中心作用$N=t+1$は$N^2=0$である。
$s=i+N/(2i)$と置くと$s^2=t$。
$X=\bar x/s$, $Y=\bar y/s$, $Z=\bar z/s$と置けば
$E$は$R_1$-加群になり、
\[
 G_1=-b(s),\qquad b(s)=s^{-3}+s^3
\]
を満たす。
(QS)により$R_1$-自己拡大は分裂するため
$G_1$は$E$上でも零。
$b(i)=0$、$b'(i)=-6$、$(s-i)^2=0$から
$s=i$、$t=-1$を得る。
分裂は$Q^\circ$-線形となり、Extが零である。
大域次元については$s^2=t$を添加すると
\[
 g=-\frac32(s^3+s^{-3}),\quad
 g^2-9=\frac{9(s^6-1)^2}{4s^6}
\]
となる。$t^3-1$を逆にした領域では
$s\mapsto g$の微分は単元で、
$R_1[(g^2-9)^{-1}]$からの有限\'etale中心基底変換となる。
有限\'etale拡大での大域次元保存と忠実平坦降下から
上界三を得る。
\end{proof}

\begin{lemma}[coherent対象への導来随伴]
\label{lem:quotient-chart-ext}
上記の仮定の下で
\[
 \mathbf R\operatorname{Hom}(\mathcal O,\mathcal T)
   \simeq V[0],\quad
 \operatorname{End}(\mathcal T)=\CC,\quad
 \operatorname{Ext}^1(\mathcal T,\mathcal T)=0 .
\]
また$\mathcal T$は大きい$\operatorname{QGr}A$内で
長さ高々三のinjective分解を持つ。
\end{lemma}
\begin{proof}
$A^\circ=Q^\circ[u,u^{-1}]$と書ける。
中心局所化から
$j^*:\operatorname{QGr}A\to\operatorname{Mod}Q^\circ$
を得る。その右随伴
$j_*M=\pi(M\otimes_{Q^\circ}A^\circ)$は、
$A^\circ$が左$Q^\circ$加群として自由なのでexactである。
ここで局所化済み加群$L=M\otimes_{Q^\circ}A^\circ$はtorsion-freeかつ飽和している。
実際、$A_{\ge n}\to L$という次数付き写像は、中心元$u^n$の像を
$u^{-n}$倍することで$A\to L$へ一意に延長する。
$V$から作る非負次数格子$N_V$に対して$L/N_V$はtorsionである
（各負次数元は十分高い正次数で非負次数へ入る）。
従って$j_*V\simeq\pi N_V=\mathcal T$である。
$j^*$もexactだから導来随伴
\[
 \mathbf R\operatorname{Hom}_{\operatorname{QGr}A}
 (F,j_*V)\simeq
 \mathbf R\operatorname{Hom}_{Q^\circ}(j^*F,V)
\]
が成立する。$j^*\mathcal O=Q^\circ$、
$j^*\mathcal T=V$を代入して前半を得る。
$\operatorname{QGr}A$は局所noetherianなGrothendieck圏で、
coherent対象間のExtは$\operatorname{qgr}A$内と一致する。
最後に$\operatorname{gldim}Q^\circ\le3$で
$V$の長さ三以下のinjective分解を取り、
exactな$j_*$で移せばよい。
\end{proof}

\begin{theorem}[構造対象をline objectのshiftに送らない自己同値]
\label{thm:quotient-spherical-counterexample}
上記で証明した$\alpha=1$の正則性とSerre関手を用いる。
(QS)は命題\ref{prop:local-two-axial}、\ref{prop:quotient-axial-rigidity}により成立する。$\mathcal T$は3-spherical対象であり、
Seidel--Thomasの球面ツイスト
\cite[Definition~2.14, Proposition~2.10]{SeidelThomas}
\[
 \mathbb T_{\mathcal T}(F)=
 \operatorname{Cone}\left(
 \mathbf R\operatorname{Hom}(\mathcal T,F)
    \otimes_\CC\mathcal T\longrightarrow F
 \right)
\]
は$\Db(\operatorname{qgr}A)$の自己同値を与える。
標準$t$構造に関して
\[
 \mathcal H^j(\mathbb T_{\mathcal T}(\mathcal O))=
 \begin{cases}
 \mathcal O,&j=0,\\
 \mathcal T^{\oplus3},&j=2,\\
 0,&j\ne0,2.
 \end{cases}
\]
従って全ての$m,r\in\mathbb Z$について
$\mathbb T_{\mathcal T}(\mathcal O)\not\simeq\mathcal O(m)[r]$。
\end{theorem}
\begin{proof}
$\mathcal T(1)\simeq\mathcal T$だから
$\mathbb S\mathcal T\simeq\mathcal T[3]$。
補題\ref{lem:quotient-chart-ext}とSerre双対から
\[
 \operatorname{Ext}^j(\mathcal T,\mathcal T)=
 \begin{cases}\CC,&j=0,3,\\0,&j\ne0,3.\end{cases}
\]
負次数はheartの二対象間のExtで消え、
高次数は$\operatorname{gldim}=3$により消える。
Serre双対性は必要な合成対の非退化性を与える。
補題の有限injective分解と
coherent対象間の有限次元Extにより
Seidel--Thomasの球面対象の条件を満たす。
その際、大きい圏$\operatorname{QGr}A$は局所noetherianな
Grothendieck圏なので十分なinjectiveを持ち、injectiveの直和もinjectiveである。
小さい圏$\operatorname{qgr}A$はnoetherian対象からなるSerre部分圏で、
全ての大きい対象はcoherent部分対象の有向和である。
従ってcoherent対象への全射にはcoherent部分対象からの全射が存在し、
Seidel--Thomasの条件(C1)--(C4)を満たす。
有限大域次元とproper性により、各coherent対象との全Extも有限次元である。
これでinjective enhancement上の評価写像とconeの構成が適用できる。
これは同論文の一般Abel圏での構成であり、
可換スキーム上の対象に限定しない。

さらにSerre双対と上記随伴から
\[
 \mathbf R\operatorname{Hom}(\mathcal T,\mathcal O)
  \simeq
  \mathbf R\operatorname{Hom}
     (\mathcal O,\mathbb S\mathcal T)^*
  \simeq V^*[-3].
\]
よってツイストの三角形は
$\mathcal T^{\oplus3}[-3]\to\mathcal O
\to\mathbb T_{\mathcal T}(\mathcal O)
\to\mathcal T^{\oplus3}[-2]$となり、
cohomology長完全列から表示が従う。
$\mathcal O$と$\mathcal T$はどちらも非零なので、
二つの異なる次数に非零cohomologyを持ち、
単一coherent対象のshiftである$\mathcal O(m)[r]$とは同型でない。
\end{proof}
\begin{remark}[標識なし復元とは異なる]
\label{rem:quotient-unmarked-reconstruction}
この自己同値は「任意の導来同値が構造対象のtwist/shiftを
保つ」という自動性を否定する。
異なる二つの代数の標識なし導来圏が同値なら
次数付き同型かどうか、また球面ツイスト等を合成すれば
標識を戻せるかどうかまでは決定しない。
従って構造対象を指定した復元定理とは矛盾しない。
\end{remark}
